%% ============================================================================
\chapter{Results}
\label{ch:results}
\addcontentsline{lot}{chapter}{Results}
%% ============================================================================

%% ---------------------------------------------------------------------------
%% SCAFFOLD. The structure, the contrasts and the reporting conventions are
%% fixed here so that the experiments have a defined target. Every numeric claim
%% is marked \PLACEHOLDER until the corresponding cluster run has completed.
%% Do NOT populate these tables from the workstation environment: see
%% Appendix A, Section A.3 for why those numbers are not comparable.
%% ---------------------------------------------------------------------------

This chapter reports the evaluation described in
Section~\ref{sec:evaluation-plan}.
Section~\ref{sec:res-rq1} establishes when the loader is the bottleneck (RQ1);
Sections~\ref{sec:res-rq2}--\ref{sec:res-rq4} address sharding, predictive staging
and prefetch-depth auto-tuning in turn.
Section~\ref{sec:res-accuracy} verifies the accuracy constraint.

All figures are reported as mean $\pm$ standard deviation over
\PLACEHOLDER{$n$} repeats, since measurements on a shared filesystem are
non-stationary and a single run is not evidence.

\section{RQ1: When does the DataLoader become the bottleneck?}
\label{sec:res-rq1}

\begin{table}[!ht]
\centering
\begin{tabular}{llccccc}
\toprule
Architecture & Loader & \Tdata{} & \Tdecode{} & \Thtod{} & \Tgpu{} & GPU util.\ (\%) \\
\midrule
SegFormer   & default & & & & & \\
SegFormer   & tuned   & & & & & \\
MA-Net      & default & & & & & \\
MA-Net      & tuned   & & & & & \\
DeepLabV3+  & default & & & & & \\
DeepLabV3+  & tuned   & & & & & \\
\bottomrule
\end{tabular}
\caption{Per-stage timing breakdown across architectures and loader
configurations. \PLACEHOLDER{populate from cluster profiling runs}}
\label{tab:rq1-breakdown}
\end{table}

The prediction under test is directional and stated in advance: SegFormer, having
the lowest \Tgpu{} of the three architectures
(Table~\ref{tab:baseline-default}), should be the configuration in which the
loader bottleneck appears first, because a model that computes quickly spends
proportionally more of its step waiting.
DeepLabV3+, the heaviest, should be the configuration in which staging buys least.
If the ordering does not hold, the framing of this thesis is wrong and should be
reported as such.

\PLACEHOLDER{narrative: does the predicted ordering hold?}

\section{RQ2: Do tar shards reduce many-small-file overhead?}
\label{sec:res-rq2}

The direct measure is filesystem operation counts, not wall-clock: sharding is
intended to convert many small metadata operations into few large sequential
reads, and that intent is verifiable independently of whether it happens to pay off
in this configuration.

\begin{table}[!ht]
\centering
\begin{tabular}{lccc}
\toprule
Method & \code{open} calls & \code{stat} calls & Bytes read \\
\midrule
1. Baseline loose      & & & \\
3. WebDataset (shards) & & & \\
5. \pads{}             & & & \\
\bottomrule
\end{tabular}
\caption{Filesystem operations issued against GPFS per epoch.
\PLACEHOLDER{populate}}
\label{tab:rq2-fsops}
\end{table}

\PLACEHOLDER{narrative: operation-count reduction, and whether it translates into
wall-clock benefit}

A negative result here is publishable and should be reported as such: if operation
counts fall substantially but throughput does not, the conclusion is that metadata
overhead was not the binding constraint on this system, which is itself an answer
to RQ1.

\section{RQ3: Does predictive staging outperform fixed streaming?}
\label{sec:res-rq3}

\begin{table}[!ht]
\centering
\begin{tabular}{lccccc}
\toprule
Method & Throughput (samples/s) & GPU util.\ (\%) & s/epoch & IoU & $\Delta$ vs.\ ref. \\
\midrule
1. Baseline \code{DataLoader}     & & & & & \\
2. Tuned \code{DataLoader}        & & & & & \\
3. WebDataset                     & & & & & \\
4. Full pre-stage (ceiling)       & & & & & \\
5. \pads{}                        & & & & & \\
\bottomrule
\end{tabular}
\caption{End-to-end comparison. Method 2 is the bar \pads{} must clear;
method 4 is the ceiling. \PLACEHOLDER{populate}}
\label{tab:rq3-endtoend}
\end{table}

Two comparisons carry the argument.
Against method~2 --- \pads{} must beat a competently hand-tuned loader, not merely
the naive default, or the contribution is a tuning guide rather than a system.
Against method~4 --- if manually pre-staging the whole corpus matches \pads{}, then
adaptivity bought nothing \emph{for this dataset}, and the honest conclusion is that
the contribution reduces to ``the dataset fits in scratch''.
The defence in that case is not that the result is otherwise, but that \pads{}
reaches the ceiling without being told scratch is the answer, and continues to
function when the corpus does not fit --- which Chapter~\ref{ch:ablation} probes
directly by constraining scratch capacity.

\PLACEHOLDER{narrative}

\section{RQ4: Can prefetch depth be auto-tuned?}
\label{sec:res-rq4}

\begin{table}[!ht]
\centering
\begin{tabular}{lcc}
\toprule
Prefetch depth $d$ & Throughput (samples/s) & Peak scratch (GB) \\
\midrule
1 (fixed)    & & \\
2 (fixed)    & & \\
4 (fixed)    & & \\
8 (fixed)    & & \\
\pads{} (auto-tuned, $d = $ \PLACEHOLDER{}) & & \\
\bottomrule
\end{tabular}
\caption{Fixed prefetch depths against the auto-tuned depth.
\PLACEHOLDER{populate}}
\label{tab:rq4-depth}
\end{table}

The claim is falsifiable in a specific way: the auto-tuned depth should land at or
near the best fixed depth \emph{without having been told it}.
If the swept optimum is materially better than the derived depth, the derivation in
Section~\ref{ch:methodology} is mis-specified and should be revised rather than
defended.

\PLACEHOLDER{narrative}

\section{Profiling overhead}

The warm-up window is a cost, and reporting throughput gains without it would
overstate the result.

\begin{table}[!ht]
\centering
\begin{tabular}{lcc}
\toprule
 & Absolute (s) & \% of total run \\
\midrule
Warm-up profiling window   & & \\
Archive construction (one-off, offline) & & \\
Net benefit after overhead & & \\
\bottomrule
\end{tabular}
\caption{Cost of adaptation. \PLACEHOLDER{populate}}
\label{tab:overhead}
\end{table}

\section{Accuracy is preserved}
\label{sec:res-accuracy}

\pads{} is only useful if it accelerates the pipeline without degrading
segmentation quality.
The relevant threshold is stated in advance: any IoU deviation exceeding the
$\pm 0.5$ percentage point standard deviation of the reference results
(Table~\ref{tab:casini-reference}) is a regression to be explained, not noise.

\begin{table}[!ht]
\centering
\begin{tabular}{lccc}
\toprule
Method & IoU (per-image, all tiles) & $\Delta$ vs.\ baseline & Within tolerance? \\
\midrule
1. Baseline & & --- & --- \\
3. WebDataset & & & \\
5. \pads{} & & & \\
\bottomrule
\end{tabular}
\caption{Segmentation quality under each data path. Reduction convention as
Section~\ref{sec:repro-metric}. \PLACEHOLDER{populate}}
\label{tab:accuracy-preserved}
\end{table}

One mechanism deserves specific attention rather than a summary statistic.
Shard-based reading approximates exact random shuffling by shuffling shard order and
maintaining a sample buffer; reduced shuffling randomness is a plausible route by
which a throughput optimisation could quietly degrade the model.
The comparison in Table~\ref{tab:accuracy-preserved} is therefore not merely a
sanity check but a test of that specific hypothesis, and
Section~\ref{sec:ablation-shuffle} varies the buffer size to probe it.

\PLACEHOLDER{narrative}
